% ЕМТ 2023, IX клас — точен LaTeX препис от официалните файлове в Klasirane.com.
% Условия: https://klasirane.com/api/competitions/EMT/All/2023/9%20%D0%BA%D0%BB/probs
% SHA-256 (условия): c0bf91a00241b5c6a510975999b6503b57f247e4c8fe910848f324fd8adfc5c0
% Решения: https://klasirane.com/api/competitions/EMT/All/2023/9%20%D0%BA%D0%BB/sol
% SHA-256 (решения): 2e9392570b447672d3432efbda76777427fed8fe62febf4dc127538ff70f6c21
% Точковите оценки, правописът, формулировките и видимите печатни особености са запазени от източника.
% В официалните файлове няма отделна чертежна фигура; не е добавяна такава.

Задача 9.1. Дадени са функциите $f(x)=|x-2|-|x-4|$ и $g(x)=|x-8|-2$. Да се пресметне лицето на фигурата с върхове, пресечните точки на графиките на функциите $f(x)$ и $g(x)$ и пресечните точки на графиката на $g(x)$ с оста $Ox$.

Решение. Графиката на $g(x)$ се състои от два лъча с общ връх в $x=8$. Разкриваме модула и лесно изчисляваме пресечните й точки с оста $Ox$ чрез уравненията $6-x=0$ и $x-10=0$ - $A(6,0)$ и $B(10,0)$.

След разкриване на модулите в $f(x)$ виждаме че в интервала $(-\infty,2)$ $f(x)=-2$, в интервала $[2,4]$ $f(x)=2x-6$ и в интервала $(4,\infty)$ $f(x)=2$.

Решаваме уравненията $f(x)=g(x)$ за всеки от трите интервала. Получаваме следните пресечни точки - $D(4,2)$ и $C(12,2)$.

От координатите следва, че фигурата $ABCD$ е трапец с основи 8 и 4 и височина 2. Следователно лицето му е $\dfrac{8+4}{2}=12$.

Оценяване. (6 точки) 2т за пресечните точки на $g(x)$ с $Ox$. 2т. за пресечните точки на $g(x)$ с $f(x)$. 1т. за това че разглежданата фигура е трапец. 1т. за довършване.

Задача 9.2. Даден е тъпоъгълен равнобедрен триъгълник $ABC$ ($AC=BC$), около който е описана окръжност с център $O$. Точка $P$ е произволна точка върху основата $AB$, такава че $AP<\dfrac12AB$. Точка $Q$ лежи на основата $AB$ и $BQ=AP$. Окръжността с диаметър $CQ$ пресича описаната около триъгълник $ABC$ окръжност за втори път в точка $E$, а правите $CE$ и $AB$ се пресичат в точка $F$. Ако $N$ е средата на $CP$ и правите $ON$ и $AB$ се пресичат в точка $D$, да се докаже че точките $O,D,C,F$ лежат на една окръжност.

Решение. Нека $T$ е среда на $CQ$ и нека означим $\angle NOC=\angle TOC=\alpha$. $OT$ е перпендикулярна на $CE$ защото $T$ е център на окръжността с диаметър $CQ$. Нека $K$ е пресечна точка на $ON$ и $CF$. Намираме $\angle OKC=\angle OKF=90-2\alpha$. От $\angle PDN=\angle PDK=90-\alpha$ намираме $\angle DFC=\alpha$ т.е. $\angle COD=\angle CFD=\alpha$. Това завършва доказателството.

Оценяване. (6 точки) 2т. за въвеждане на $\angle TOC$; 1т. за перпендикулярността; 1т. за $\angle OKC=\angle OKF=90-2\alpha$; 1т. за $\angle DFC=\alpha$; 1т. за довършване.

Задача 9.3. В къщата на богатата лейди Гилмор се случила кражба на една от най-скъпите й ценности: нейната перлена огърлица. Задачата за разплитането на мистерията паднала на плещите на инспектор Гудинаф. Той разполагал със следната информация: в деня на кражбата, в стаята с огърлицата били влизали 7 от слугите на лейди Гилмор, които ще наричаме $A,B,C,D,E,F,G$ поради конфиденциалност на разследването. Всеки от тях твърди, че е присъствал в стаята само веднъж за неопределен период от време. Освен това $A$ твърди, че е срещал $B,C,F,G$ в стаята; $B$ твърди, че е срещал $A,C,D,E,F$; $C$ твърди, че е срещал $A,B,E$; $E$ твърди, че е срещал $B,C,F$; $F$ твърди, че е срещал $A,B,D,E$; $G$ твърди, че е срещал $A,D$ и $D$ твърди, че е срещал $B,F,G$. Инспектор Гудинаф заключил, че точно един от слугите лъже. Кой е той?

Решение. Първо ще докажем следната лема.

Лема. Нека $X,Y,Z$ и $T$ са четирима от слугите. Ако е известно, че двойките $X,Y$; $Y,Z$; $Z,T$ и $T,X$ са били заедно в стаята в даден момент, то някоя от двойките $X,Z$ и $Y,T$ също са се засекли.

Доказателство на Лема. Нека без ограничение на общността допуснем, че $Y$ и $T$ не са били заедно в стаята и $Y$ си е тръгнал от стаята преди $T$ (останалите случаи са аналогични). Тогава, $X$ и $Z$ са стояли в стаята заедно в периода между напускането на $Y$ и пристигането на $T$.

Да забележим, че $A,C,E,F$ удовлетворяват условието на лемата, но никои от $A,E$ и $C,F$ не са се засякли. Същото важи за $A,B,D,G$. Единствен общ елемент на тези двойки е $A$. Остава да се уверим, че е възможно всички останали двойки да са се срещнали, както твърдят, влизайки точно по веднъж. Това е възможно при следната последователност от влизания и излизания: влиза $G$, влиза $D$, излиза $G$, влиза $B$, влиза $F$, излиза $D$, влиза $E$, излиза $F$, влиза $C$, излиза $B$, излиза $E$, излиза $C$.

Оценяване. (7 точки) 2т за твърдението на лемата; 2т за правилно доказателство на лемата; по 1т за всяка четворка, изобличаваща $A$; 1т за пример, че всички освен $A$ казват истината.

Задача 9.4. Нека $p$ и $q$ са взаимнопрости цели числа и $\left|\dfrac pq\right|\le1$. Да се определи за кои стойности на $p$ и $q$ съществува представяне от вида
$$
\frac pq=\cfrac{1}{b_1+\cfrac{1}{b_2+\cfrac{1}{b_3+\cdots}}}
$$
за краен брой четни числа $b_1,b_2,\ldots,b_n$.

Бележка: Това представяне се нарича верижна дроб и може да бъде означавано и като $[b_1,b_2,\ldots,b_n]$.

Решение. Ще докажем че това е възможно точно когато едно от двете числа $p$ и $q$ е четно.

1) Необходимост. Ако всички $b_i$ са четни, ще докажем че $pq$ е четно. Ще използваме индукция по $n$. За $n=1$ това е очевидно.

Нека направим следното представяне:
$$
\frac pq=[b_1,\ldots,b_n]=\frac{1}{b_1+[b_2,\ldots,b_n]}=\frac{1}{b_1+\frac{p'}{q'}}=\frac{q'}{b_1q'+p'}
$$

Тук $b_1$ е четно и по индукционно допускане точно едно от $p'$ и $q'$ също е четно. Лесно се вижда това завършва доказателството на тази посока.

2) Достатъчност. Ще докажем, че съществуват такива четни числа $[b_1,\ldots,b_n]$ при $pq$ четно. Отново ще ползваме индукция, но този път по $|q|$.

За $|q|=2$ твърдението отново е очевидно.

Нека разгледаме числата $\left\lfloor\dfrac qp\right\rfloor$ и $\left\lfloor\dfrac qp\right\rfloor+1$. Едно от тях е четно. Да го наречем $b$ и да отбележим, че не може да е 0. Числото $b-\dfrac qp$ е по-малко от единица и може да бъде записано като несъкратима дроб $\dfrac{p'}{q'}=b-\dfrac qp$ или алтернативно $\dfrac pq=\dfrac{1}{b+\frac{p'}{q'}}$, като $|q|>|q'|$. Сега аналогично на предния случай можем да докажем че $p'q'$ е четно и $\dfrac{p'}{q'}$ има необходимото представяне според индукционното допускане. Но тогава и $\dfrac pq$ има такова. Това завършва доказателството.

Оценяване. (7 точки) 1т за правилен отговор. 3т. за всяка посока на доказателството.

Бележка: Това твърдение, както и самите верижни дроби, имат приложение в теория на възлите. Също така е доказано и единственост на представянето, макар че това свойство не се иска в задачата.
